A (central) \textit{hyperplane arrangement} $ \mathcal{A} $ over a field $ \mathbb{K} $ is a finite collection of subspaces of codimension $ 1 $ in a finite dimensional vector space $ \mathbb{K}^{\ell} $.
A standard reference for arrangements is~\cite{orlik1992arrangements}.
Let $ S $ denote the polynomial algebra $ \mathbb{K}[x_{1}, \dots, x_{\ell}] $, where $ (x_{1}, \dots, x_{\ell}) $ is a basis for the dual space $ \big(\mathbb{K}^{\ell}\big) ^{\ast} $.
Let $ \Der(S) $ denote the \textit{module of derivations} of $ S $, that is,
\begin{gather*}
\Der(S) \coloneqq \{ \theta \colon S \to S \,|\, \theta \text{ is $ S $-linear and } \theta(fg) = f\theta(g) + \theta(f)g \text{ for any } f,g \in S \}.
\end{gather*}
The \textit{module of logarithmic derivations} $ D(\mathcal{A}) $ is defined by
\begin{gather*}
D(\mathcal{A}) \coloneqq \{ \theta \in \Der(S) \,|\, \theta(\alpha_{H}) \in \alpha_{H}S \text{ for all } H \in \mathcal{A} \},
\end{gather*}
where $ \alpha_{H} $ is a linear form such that $ \ker(\alpha_{H}) = H $.

\begin{Definition}
An arrangement $ \mathcal{A} $ is called \textit{free} if $ D(\mathcal{A}) $ is a free $ S $-module.
\end{Definition}

Although the definition of free arrangements is algebraic, Terao's celebrated factorization theorem \cite[Main Theorem]{terao1981generalized-im} shows a solid relation between algebra, combinatorics, and topology of arrangements.
Terao's conjecture asserts that the freeness of an arrangement is determined by its combinatorial property and it is still widely open.
